\section{Reproducibility, evidence boundaries, and integration}
\label{app:reproducibility}

\subsection{What the package contains}

The archive contains the article source, its component sections, the
compiled PDF, vector figures with plotting sources, exact verification
programs, numerical diagnostics, and machine-readable receipts. A
provenance manifest records the pinned source commit and the SHA-256
digests of the consulted manuscript files. A separate claim ledger maps
results to their proofs, their scripts, and their status. The two narrow
text corrections are provided as an ordinary unified patch, together with
integration notes keyed to existing manuscript labels.

The article is a self-contained continuation, not a silent replacement
of the 228-page source manuscript. Each contribution is isolated in a
section file so that a repository maintainer can move it to the relevant
chapter. Numbered labels have distinct prefixes for the relation-lattice,
mixed-family, CM, and classical-zero contributions. The common zero
preliminaries can be omitted when integrating into the existing zero
chapter, provided the references are redirected to its established
theorems.

\subsection{Exact computations and analytic certificates}

\begin{center}
\small
\begin{tabularx}{\textwidth}{@{}>{\raggedright\arraybackslash}p{0.31\textwidth}>{\raggedright\arraybackslash}X>{\raggedright\arraybackslash}p{0.21\textwidth}@{}}
\toprule
Program & Verified object & Logical status\\
\midrule
\texttt{verify\_rank.py} & Integer rows, rational ranks, primitive kernel
witnesses, and Smith forms in the specified finite range. & Exact
algebra; checks the general proof.\\[3pt]
\texttt{verify\_mixed\_family.py} & A finite integer row certificate for
twice the target at even weights; independent integral evaluation. &
Exact rows; quadrature is diagnostic.\\[3pt]
\texttt{cm\_norms.py} & Reduced forms, interval-enclosed singular moduli,
unique integer class-polynomial coefficients, exact factorizations. &
Certified computation using CM integrality.\\[3pt]
\path{cm_genus.py} & Quadratic-field factors, certified genus labels,
and exact cubic and square identities for the ratios. & Exact and interval certificate.\\[3pt]
{\footnotesize\texttt{modular\_weight\_polynomials.py}} & Exact $F_w$ and coefficient
identities through the weight-$12w$ valence bound, $4\le w\le24$. & Exact modular-form certificate.\\[3pt]
\path{certify_n5_zeros.py} & Rational interval signs at all eleven
root brackets and all eight needed critical values; three index-three
endpoint signs. & Finite analytic certificate.\\[3pt]
\path{certify_lerch_splitting.py} & 120 signs of integer polynomials
at $\log2$, for $2\le n\le16$. & Exact rational enclosures.\\[3pt]
\texttt{lerch\_fold\_diagnostic.py} & A high-precision approximate
solution of $D=\partial_aD=0$. & Numerical diagnostic only.\\[3pt]
\texttt{explore\_n6\_n7.py} & Finite scans and root refinements supporting
Conjecture~\ref{conj:n6n7}. & Numerical evidence only.\\
\bottomrule
\end{tabularx}
\end{center}

The zero certificate uses only Python's integer and rational arithmetic;
its logarithm bounds and Euler--Maclaurin remainder are proved in the
text. It does not call a numerical Stieltjes-constant routine. The matrix
certificate uses exact SymPy arithmetic and an independently constructed
matrix. The mixed-family certificate works directly with product rows,
and the all-weight proof supplies the uniform row combination.

The six radical genus evaluations also have an independent checker,
\path{verify_genus_arithmetic.py}, which performs arithmetic as rational
pairs in the two quadratic fields and checks the positive-root choices.
It imports neither SymPy nor the CM verifier.

The CM coefficient certificate uses the directed interval implementation
in \texttt{mpmath.iv}. Its proof contract is explicit: the elementary
interval functions enclose their true values, the rational tail bound is
added before division or multiplication, and each resulting coefficient
interval has a unique possible integer. Classical CM integrality supplies
the fact that the true coefficient is an integer. This is a conventional
computer-assisted proof, with the interval implementation as part of its
trusted computational base. No claim of proof-assistant formalization is
made.

\subsection{Replay and interpretation}

From the extracted archive, the principal commands are
\begin{verbatim}
python3 -m pip install -r requirements.txt
python3 code/verify_all.py
latexmk -pdf -interaction=nonstopmode -halt-on-error article.tex
\end{verbatim}
The replay driver runs each principal verifier in a fresh output
directory, preserving the delivered receipts for comparison. Numerical
diagnostics can be run separately with the documented flag. They are
excluded from the proof-only interpretation even when they agree to the
full working precision.

For the rank replay, the delivered run covers weights $3$ through $13$
and both Smith forms through weight $12$. It checks the $92\times23$
weight-five matrix and the exact pairing $7$ of the short witness with
the target. For the mixed family, the finite row certificate is checked
at all even weights $4$ through $20$, and the independent integral at
weights $2,4,6,8,10$. The general claims are proved in the article;
these finite ranges are validation ranges, not the asserted scope of
the theorems.

The article's PDF was compiled with all cross-references resolved and
rendered for visual inspection. Figure labels, tables, long formulas,
and page breaks were checked. The archive includes a checksum manifest
so that later editorial changes can be distinguished from the delivered
research version.

\subsection{Suggested integration order}

First apply the two local corrections. Next insert the relation-lattice
theorem after the formal product system and replace the old odd-weight
rank conjecture by a reference to the theorem. Place the mixed-family
identity next to the existing Gaussian and signed-color evaluations.
In the CM chapter, upgrade the product identities only after inserting
the period normalization, certified class data, and phase conventions.
In the zero chapter, put the fifth-index transition after the small-index
certificates, and the small-parameter and bifurcation results after the
definition of the Lerch family. Retain the diagnostic labels on the
fold coordinates and the sixth- and seventh-index conjectures.

This integration order preserves the distinction between a theorem,
an exactly checkable finite lemma, an independently verified numerical
identity, and an open conjecture throughout the combined manuscript.
